\section{Interactividad}

\textit{Matplotlib} no solo genera imágenes estáticas; posee la capacidad de interactuar con el sistema operativo para actualizar una ventana de visualización en tiempo real. Esto es fundamental para monitorear datos que cambian constantemente o para realizar demostraciones dinámicas.

\subsection{El Modo Interactivo (\textit{plt.ion})}

Por defecto, \textit{Matplotlib} trabaja en modo bloqueante: el script se detiene hasta que usted cierra la ventana del gráfico. Al activar el Interactive On (\texttt{plt.ion()}), permitimos que el código siga ejecutándose mientras la ventana permanece abierta y se actualiza.

\subsubsection{Componentes del bucle interactivo}

Para que una animación funcione correctamente sin colgar el sistema, necesitamos coordinar tres funciones del objeto \texttt{canvas} (el lienzo interactivo):

\begin{itemize}
\item \texttt{fig.canvas.draw\_idle()}: Solicita al programa que redibuje la figura en el próximo momento que el procesador esté libre. Es más eficiente que un redibujado forzado.
\item \texttt{fig.canvas.flush\_events()}: Limpia la cola de eventos del sistema operativo (clics, movimientos de mouse), permitiendo que la ventana responda y no aparezca como ``No responde''.
\item \texttt{line.set\_ydata(y)}: En lugar de borrar y volver a crear el gráfico (lo cual es lento), solo actualizamos los datos de la línea existente. Esto garantiza una animación fluida.
\end{itemize}

Veamos un ejemplo sencillo
\begin{pycode}{Simulación de Oscilación}
import numpy as np
import matplotlib.pyplot as plt
import time

# Activamos la interactividad
plt.ion()

# Creamos la figura
fig, ax = plt.subplots(layout="constrained")

# Datos
x = np.linspace(0, 10, 100)
line, = ax.plot(x, np.sin(x))

# Realizamos la simulación
for i in range(50):
    nuevo_y = np.sin(x + i/10.0) # Desplazamos la fase
    line.set_ydata(nuevo_y)      # Actualizamos solo los datos
    fig.canvas.draw_idle()
    fig.canvas.flush_events()
    time.sleep(0.05)             # Velocidad de animación

# Desactivamos la interactividad
plt.ioff()
plt.show()
\end{pycode}

\begin{nota}
Para que esto funcione fuera de un Notebook, su instalación de \textit{Python} debe tener acceso a un kit de herramientas de interfaz gráfica. Los más comunes son \textit{TkAgg} (viene con Python) o \textit{QtAgg} (más moderno y rápido, requiere instalar \texttt{PyQt5}).
\end{nota}


\subsection{TIpos de interactividad}

\subsubsection{}
